\textbf{COLUNI 2019 -- Questão 37}\par\medskip Veja o Cartaz do filme “Olga” (2004) e leia a sua sinopse:

% COLUNI-AUTO-FIGURES-BEGIN
\begin{figure}[htbp]
\centering
\includegraphics[width=0.95\linewidth]{static/simulados/coluni/figuras/2019/questao-37/figura-1.png}
\end{figure}
% COLUNI-AUTO-FIGURES-END

\par

\textbf{SINOPSE}

\par

O filme conta a história da judia alemã Olga Benário (Camila Morgado). Militante comunista desde jovem, Olga é perseguida pela polícia e foge para Moscou, onde faz treinamento militar. Em Moscou, Olga é encarregada de acompanhar Luís Carlos Prestes (Caco Ciocler) ao Brasil para liderar a Intentona Comunista de 1935, se apaixonando por ele na viagem. Com o fracasso da revolução, Olga é presa com Prestes. Grávida de 7 meses é deportada pelo governo Vargas para a Alemanha Nazista e tem sua filha Anita Leocádia na prisão. Afastada da filha, Olga é enviada para um campo de concentração de Ravensbrück.

\par

(Disponível em: <https://cinemahistoriaeducacao.wordpress.com/cinema-e-historia/brasil-republica/olga/>. Acesso em: 3 out. 2018.)

\par

A história de Olga Benário Prestes mostrada no filme ocorre durante a Era Vargas (1930-1945) em que Getúlio Vargas esteve no comando político do Brasil como Presidente da República por 15 Anos. A partir de 1937, o Regime passou a se chamar “Estado Novo”.

\par

Sobre esse período e considerando a sinopse do filme “Olga”, assinale a afirmativa INCORRETA:

\begin{enumerate}
\item O filme não retrata a realidade política da época, pois o Governo Vargas ficou do lado dos Estados Unidos
        e seus aliados contra as potências do Eixo, na defesa da liberdade e da democracia, durante a Segunda
        Guerra Mundial (1939-1945).
\item Durante o período do Estado Novo, o Departamento de Imprensa e Propaganda (DIP) controlava a
        produção cultural, a imprensa e procurava divulgar a imagem de Getúlio Vargas junto à classe
        trabalhadora.
\item No período do Estado Novo, os opositores do regime (intelectuais, simpatizantes ou membros do Partido
        Comunista do Brasil) foram perseguidos e mantidos como presos políticos, assim como foram os casos de
        Olga Benário e seu companheiro, Luiz Carlos Prestes.
\item O Governo Vargas estimulava uma produção cultural ufanista de cunho nacionalista que valorizava as
        belezas naturais do país e de cartilhas escolares que procuravam doutrinar as crianças na ideologia do
        Estado Novo, de reverência ao Presidente.\end{enumerate}
